\begin{helper}
Fix a positive integer $\Delta$.  For a fixed independent triple
$\{a,b,c\}$, write
\[
\ell_{ab}=\frac{C_{ab}}{\Delta},\qquad
\ell_{ac}=\frac{C_{ac}}{\Delta},\qquad
\ell_{bc}=\frac{C_{bc}}{\Delta},
\]
where the three $C$'s are the common-neighbor pattern counts supplied by
\noderef{SamplingTriplePatternCountingBridge}.  Suppose the paired
ordered-chamber middle exponents and lower exponents satisfy
\[
N_{ab}+C_{ab}=\Delta,\quad
N_{ac}+C_{ac}=\Delta,\quad
N_{bc}+C_{bc}=\Delta,
\]
and
\[
L_{ab}+C_{ac}+C_{bc}=\Delta+n_{abc},\quad
L_{ac}+C_{ab}+C_{bc}=\Delta+n_{abc},\quad
L_{bc}+C_{ab}+C_{ac}=\Delta+n_{abc}.
\]
Put
\[
x=1-\ell_{ab},\qquad y=1-\ell_{ac},\qquad z=1-\ell_{bc},
\qquad \tau=\frac{n_{abc}}{\Delta},
\]
and
\[
\rho_x=\frac{L_{ab}}{\Delta},\qquad
\rho_y=\frac{L_{ac}}{\Delta},\qquad
\rho_z=\frac{L_{bc}}{\Delta}.
\]
If $x,y,z>0$, then the hypotheses required by the subsequent full-chamber
integral estimate hold for the three paired chambers:
\[
\frac{N_{ab}}{\Delta}=x,\quad
\frac{L_{ab}}{\Delta}=\rho_x,\quad
\frac{N_{ac}}{\Delta}=y,\quad
\frac{L_{ac}}{\Delta}=\rho_y,\quad
\frac{N_{bc}}{\Delta}=z,\quad
\frac{L_{bc}}{\Delta}=\rho_z,
\]
$\rho_x,\rho_y,\rho_z,\tau\ge0$, and
\[
1+x+\rho_x=1+y+\rho_y=1+z+\rho_z=x+y+z+\tau .
\]
\end{helper}
\begin{proof}
The three middle-ratio identities are immediate from the complement
equations.  For example, from $N_{ab}+C_{ab}=\Delta$ and
$\ell_{ab}=C_{ab}/\Delta$,
\[
\frac{N_{ab}}{\Delta}
=\frac{\Delta-C_{ab}}{\Delta}
=1-\frac{C_{ab}}{\Delta}
=1-\ell_{ab}=x.
\]
The same calculation gives $N_{ac}/\Delta=y$ and
$N_{bc}/\Delta=z$.  The lower-ratio identities are definitions of
$\rho_x,\rho_y,\rho_z$, and the nonnegativity of these three parameters
and of $\tau$ follows from nonnegativity of the natural-number counts and
positivity of $\Delta$.

It remains only to check the common denominator.  Using
$L_{ab}+C_{ac}+C_{bc}=\Delta+n_{abc}$,
\[
\rho_x
=\frac{L_{ab}}{\Delta}
=1-\frac{C_{ac}}{\Delta}-\frac{C_{bc}}{\Delta}
  +\frac{n_{abc}}{\Delta}
=1-\ell_{ac}-\ell_{bc}+\tau .
\]
Therefore
\[
1+x+\rho_x
=1+(1-\ell_{ab})+(1-\ell_{ac}-\ell_{bc}+\tau)
=x+y+z+\tau .
\]
The two remaining lower-count equations give, in the same way,
$1+y+\rho_y=x+y+z+\tau$ and
$1+z+\rho_z=x+y+z+\tau$.  Thus the six ordered chambers have been grouped
into the three opposite pairs with exactly the normalized parameters used by
the subsequent full-chamber integral estimate.
\end{proof}
